\section{Basic tools for the analysis of functions over finite fields}

Throughout this section, let $q=p^s$ for a prime $p$ and $s\in\bbN$ with $s\ge 1$.
All vectors are in $\bbF_q^n$, and all expectations are uniform over the indicated finite sets.
Let $\omega_p = e^{2\pi i/p}$.

\begin{definition}[Field trace]
\label{def:field-trace}
\lean{Algebra.trace}
\leanok
The field trace $\Tr: \bbF_q \to \bbF_p$ is
\[
    \Tr(z) \coloneqq \sum_{i=0}^{s-1} z^{p^i}.
\]
\end{definition}

\begin{definition}[Base additive character]
\label{def:base-character}
\lean{BlrPcp.baseChar}
\uses{def:field-trace}
\leanok
The base additive character $\psi:\bbF_q\to\bbC$ is
\[
    \psi(z) \coloneqq \omega_p^{\Tr(z)}.
\]
\end{definition}

\begin{definition}[Additive characters]
\label{def:additive-character}
\lean{BlrPcp.charFn}
\uses{def:base-character}
\leanok
For $\alpha\in\bbF_q^n$, define the additive character
$\chi_\alpha: \bbF_q^n\to\bbC$ by
\[
    \chi_\alpha(x) \coloneqq \psi(\inner{\alpha}{x})
    = \omega_p^{\Tr(\inner{\alpha}{x})},
\]
where $\inner{\alpha}{x}=\sum_{i=1}^n \alpha_i x_i\in\bbF_q$.
\end{definition}

\begin{definition}[Expectation]
\label{def:expectation}
\lean{BlrPcp.expectation}
\leanok
For a function $h:\bbF_q^n\to\bbC$, define its uniform expectation by
\[
    \bbE_{x\in\bbF_q^n}[h(x)]
    \coloneqq \frac{1}{q^n}\sum_{x\in\bbF_q^n} h(x).
\]
\end{definition}

\begin{definition}[Inner product]
\label{def:function-inner-product}
\lean{BlrPcp.fnInner}
\uses{def:expectation}
\leanok
For functions $h_1,h_2:\bbF_q^n\to\bbC$, define
\[
    \inner{h_1}{h_2}
    \coloneqq \bbE_{x\in\bbF_q^n}\left[h_1(x)\overline{h_2(x)}\right].
\]
\end{definition}

\begin{definition}[$L_2$ norm]
\label{def:function-l2-norm}
\lean{BlrPcp.fnNorm}
\uses{def:expectation,def:function-inner-product}
\leanok
For a function $h:\bbF_q^n\to\bbC$, define its $L_2$ norm by
\[
    \|h\|_2
    \coloneqq
    \left(\bbE_{x\in\bbF_q^n}\left[\abs{h(x)}^2\right]\right)^{1/2}
    =
    \sqrt{\inner{h}{h}}.
\]
\end{definition}

\begin{definition}[Fourier coefficient]
\label{def:fourier-coefficient}
\lean{BlrPcp.fourierCoeff}
\uses{def:function-inner-product,def:additive-character}
\leanok
For $h:\bbF_q^n\to\bbC$ and $\alpha\in\bbF_q^n$, define
\[
    \widehat h(\alpha)
    \coloneqq \inner{h}{\chi_\alpha}
    = \bbE_{x\in\bbF_q^n}\left[h(x)\overline{\chi_\alpha(x)}\right].
\]
\end{definition}

\begin{lemma}[Characters are orthonormal]
\label{lem:character-orthonormal}
\lean{BlrPcp.characters_orthonormal}
\uses{def:additive-character,def:function-inner-product}
\leanok
For every $\alpha,\beta\in\bbF_q^n$,
\[
    \inner{\chi_\alpha}{\chi_\beta}
    =
    \begin{cases}
        1 & \text{if } \alpha=\beta,\\
        0 & \text{if } \alpha\ne\beta.
    \end{cases}
\]
\end{lemma}

\begin{proof}\leanok
First observe that for every $\alpha,\beta\in\bbF_q^n$ and every $x\in\bbF_q^n$,
\[
    \chi_\alpha(x)\overline{\chi_\beta(x)}
    =
    \omega_p^{\Tr(\inner{\alpha}{x})}
    \omega_p^{-\Tr(\inner{\beta}{x})}
    =
    \omega_p^{\Tr(\inner{\alpha-\beta}{x})}
    =
    \chi_{\alpha-\beta}(x),
\]
where we used the $\bbF_p$-linearity of the trace map.

If $\alpha=\beta$, then $\alpha-\beta=0$, so $\chi_{\alpha-\beta}(x)=1$ for every
$x\in\bbF_q^n$. Hence
\[
    \bbE_{x\in\bbF_q^n}\left[\chi_\alpha(x)\overline{\chi_\beta(x)}\right]
    =
    \bbE_{x\in\bbF_q^n}[1]
    =
    1.
\]

Now suppose $\alpha\ne\beta$. Let $\gamma \coloneqq \alpha-\beta$. Then $\gamma\ne 0$, so there exists
an index $j$ such that $\gamma_j\ne 0$. Fix all coordinates of $x$ except $x_j$. For the fixed
remaining coordinates, the map
\[
    x_j\mapsto \inner{\gamma}{x}
\]
has the form
\[
    x_j\mapsto \gamma_j x_j + c
\]
for some constant $c\in\bbF_q$. Since $\gamma_j\ne 0$, multiplication by $\gamma_j$ is a
bijection on $\bbF_q$, and therefore the affine map $x_j\mapsto \gamma_jx_j+c$ is also a
bijection on $\bbF_q$. Thus, as $x_j$ ranges uniformly over $\bbF_q$, the value
$\inner{\gamma}{x}$ also ranges uniformly over $\bbF_q$. Averaging first over $x_j$ and then
over the remaining coordinates gives
\[
    \bbE_{x\in\bbF_q^n}\left[\chi_\gamma(x)\right]
    =
    \bbE_{t\in\bbF_q}\left[\omega_p^{\Tr(t)}\right].
\]

It remains to show that this last average is zero. The trace map
$\Tr:\bbF_q\to\bbF_p$ is a nonzero $\bbF_p$-linear map, so every fiber of $\Tr$ has the same
cardinality. Since $\abs{\bbF_q}=q$ and $\abs{\bbF_p}=p$, each fiber has cardinality $q/p$.
Therefore
\begin{align*}
    \bbE_{t\in\bbF_q}\left[\omega_p^{\Tr(t)}\right]
    &=
    \frac{1}{q}\sum_{t\in\bbF_q}\omega_p^{\Tr(t)} \\
    &=
    \frac{1}{q}\sum_{r\in\bbF_p}
        \sum_{\substack{t\in\bbF_q\\ \Tr(t)=r}}
        \omega_p^r \\
    &=
    \frac{1}{q}\sum_{r\in\bbF_p}\frac{q}{p}\omega_p^r \\
    &=
    \frac{1}{p}\sum_{r\in\bbF_p}\omega_p^r \\
    &=0,
\end{align*}
because the sum of all $p$-th roots of unity is zero. Hence, if $\alpha\ne\beta$, then
\[
    \bbE_{x\in\bbF_q^n}\left[\chi_\alpha(x)\overline{\chi_\beta(x)}\right]
    =
    0.
\]

Combining the two cases proves the stated orthogonality relation. The equivalent formulation
follows by taking $\beta=0$.
\end{proof}

\begin{lemma}[Fourier inversion]
\label{lem:fourier-inversion}
\lean{BlrPcp.fourier_inversion}
\uses{def:fourier-coefficient,def:additive-character}
\leanok
For every function $h:\bbF_q^n\to\bbC$ we have
\[
    h(x)=\sum_{\alpha\in\bbF_q^n}\widehat h(\alpha)\chi_\alpha(x)
    \quad\text{for every }x\in\bbF_q^n.
\]
\end{lemma}

\begin{proof}\uses{lem:character-orthonormal, def:function-inner-product}\leanok
By \cref{lem:character-orthonormal}, the additive characters
\[
    \{\chi_\alpha : \alpha \in \mathbb F_q^n\}
\]
are orthonormal with respect to the inner product
\[
    \langle h_1,h_2\rangle
    =
    \mathbb E_{x\in \mathbb F_q^n}
    \left[h_1(x)\overline{h_2(x)}\right].
\]
Moreover, these characters span the vector space of all complex-valued
functions on \(\mathbb F_q^n\). Indeed, there are exactly \(q^n\) characters,
and the space of functions \(\mathbb F_q^n\to\mathbb C\) has dimension \(q^n\).
Since the characters are nonzero and pairwise orthogonal, they are linearly
independent. Hence they form an orthonormal basis.

Therefore every function \(h:\mathbb F_q^n\to\mathbb C\) has a unique expansion
in this basis:
\[
    h
    =
    \sum_{\alpha\in\mathbb F_q^n}
        c_\alpha \chi_\alpha
\]
for some coefficients \(c_\alpha\in\mathbb C\). Taking the inner product of
both sides with \(\chi_\beta\), we obtain
\[
    \langle h,\chi_\beta\rangle
    =
    \sum_{\alpha\in\mathbb F_q^n}
        c_\alpha \langle \chi_\alpha,\chi_\beta\rangle .
\]
By orthonormality, all terms in the sum vanish except the term
\(\alpha=\beta\), and \(\langle \chi_\beta,\chi_\beta\rangle=1\). Hence
\[
    \langle h,\chi_\beta\rangle = c_\beta .
\]
By definition of the Fourier coefficient,
\[
    \widehat h(\beta)=\langle h,\chi_\beta\rangle .
\]
Thus \(c_\beta=\widehat h(\beta)\) for every
\(\beta\in\mathbb F_q^n\). Therefore
\[
    h
    =
    \sum_{\alpha\in\mathbb F_q^n}
        \widehat h(\alpha)\chi_\alpha .
\]
Evaluating this identity at \(x\in\mathbb F_q^n\) gives the Fourier inversion
formula
\[
    h(x)
    =
    \sum_{\alpha\in\mathbb F_q^n}
        \widehat h(\alpha)\chi_\alpha(x).
\]
\end{proof}

\begin{lemma}[Plancherel identity]
\label{lem:fourier-plancherel}
\lean{BlrPcp.fourier_plancherel}
\uses{def:function-inner-product,def:fourier-coefficient}
\leanok
For every pair of functions $f,g:\bbF_q^n\to\bbC$,
\[
    \sum_{\alpha\in\bbF_q^n}
        \widehat f(\alpha)\overline{\widehat g(\alpha)}
    =
    \inner{f}{g}.
\]
\end{lemma}

\begin{proof}\uses{lem:fourier-inversion}\leanok
By \cref{lem:fourier-inversion},
\[
    f=\sum_{\alpha\in\bbF_q^n}\widehat f(\alpha)\chi_\alpha .
\]
Taking the inner product with \(g\) and using linearity in the first argument gives
\[
    \inner{f}{g}
    =
    \sum_{\alpha\in\bbF_q^n}
        \widehat f(\alpha)\inner{\chi_\alpha}{g}.
\]
Since
\[
    \inner{\chi_\alpha}{g}
    =
    \overline{\inner{g}{\chi_\alpha}}
    =
    \overline{\widehat g(\alpha)},
\]
we obtain
\[
    \inner{f}{g}
    =
    \sum_{\alpha\in\bbF_q^n}
        \widehat f(\alpha)\overline{\widehat g(\alpha)}.
\]
\end{proof}

\begin{lemma}[Parseval identity]
\label{lem:parseval-identity}
\lean{BlrPcp.parseval_identity}
\uses{def:fourier-coefficient,def:function-l2-norm}
\leanok
For every function $h:\bbF_q^n\to\bbC$,
\[
    \sum_{\alpha\in\bbF_q^n}\abs{\widehat h(\alpha)}^2
    =
    \|h\|_2^2.
\]
\end{lemma}

\begin{proof}\uses{lem:fourier-plancherel,def:function-inner-product}\leanok
Apply \cref{lem:fourier-plancherel} with \(f=g=h\):
\[
    \sum_{\alpha\in\bbF_q^n}
        \widehat h(\alpha)\overline{\widehat h(\alpha)}
    =
    \inner{h}{h}.
\]
The left-hand side is
\[
    \sum_{\alpha\in\bbF_q^n}\abs{\widehat h(\alpha)}^2,
\]
and by the definitions of the inner product and the $L_2$ norm, the right-hand side is
\[
    \inner{h}{h}
    =
    \bbE_{x\in\bbF_q^n}\left[\abs{h(x)}^2\right]
    =
    \|h\|_2^2.
\]
\end{proof}

\begin{definition}[Nonzero field elements]
\label{def:nonzero-field-elements}
\lean{BlrPcp.nonzeroF}
\leanok
Write
\[
    \bbF_q^\times \coloneqq \{c\in\bbF_q : c\ne 0\}.
\]
\end{definition}

\begin{definition}[Fourier expansion of a finite-field-valued function]
\label{def:fourier-expansion-field-valued}
\lean{BlrPcp.phaseLift}
\uses{def:base-character,def:nonzero-field-elements}
\leanok
Let $f:\bbF_q^n\to\bbF_q$. For each $c\in\bbF_q^\times$, define the $c$-phase lift
\[
    \fexp{\ff}{c}(x) \coloneqq \omega_p^{\Tr(c f(x))}.
\]
The Fourier expansion of $f$ is the set $\{\varphi_c\}_{c\in \bbF_q^\times}$.
\end{definition}

\begin{lemma}[Character-sum indicator]
\label{lem:character-sum-indicator}
\lean{BlrPcp.character_sum_indicator_eq}
\uses{def:base-character}
\leanok
For every $u,v\in\bbF_q$,
\[
    \ind[u=v]
    = \frac{1}{q}\sum_{c\in\bbF_q}\omega_p^{\Tr(c(u-v))}.
\]
\end{lemma}
\begin{proof}\leanok
For each $z \in \bbF_q$ we define
\[
    \psi(z) \coloneqq \omega_p^{\Tr(z)}
\]
be the standard additive character of \(\bbF_q\). We prove that, for every
\(t\in \bbF_q\),
\[
    \ind[t=0]
    =
    \frac{1}{q}\sum_{c\in\bbF_q}\psi(ct).
\]
The desired statement follows by taking \(t=u-v\).

If \(t=0\), then \(ct=0\) for every \(c\in\bbF_q\), so
\[
    \frac{1}{q}\sum_{c\in\bbF_q}\psi(ct)
    =
    \frac{1}{q}\sum_{c\in\bbF_q}1
    =
    1.
\]

Now suppose \(t\neq 0\). Then multiplication by \(t\) is a bijection
\(\bbF_q\to\bbF_q\), so
\[
    \sum_{c\in\bbF_q}\psi(ct)
    =
    \sum_{z\in\bbF_q}\psi(z).
\]
We claim that the latter sum is \(0\). Let
\[
    S \coloneqq \sum_{z\in\bbF_q}\psi(z).
\]
Since \(\Tr:\bbF_q\to\bbF_p\) is a nonzero \(\bbF_p\)-linear map, it is
surjective. Hence there exists \(a\in\bbF_q\) such that \(\Tr(a)=1\).
Using the bijection \(z\mapsto z+a\), we get
\[
    S
    =
    \sum_{z\in\bbF_q}\psi(z+a)
    =
    \sum_{z\in\bbF_q}\omega_p^{\Tr(z+a)}
    =
    \sum_{z\in\bbF_q}\omega_p^{\Tr(z)}\omega_p^{\Tr(a)}
    =
    \omega_p S.
\]
Since \(\omega_p\neq 1\), this implies \(S=0\). Therefore, when
\(t\neq 0\),
\[
    \frac{1}{q}\sum_{c\in\bbF_q}\psi(ct)=0.
\]

Combining the two cases gives
\[
    \ind[t=0]
    =
    \frac{1}{q}\sum_{c\in\bbF_q}\omega_p^{\Tr(ct)}.
\]
Substituting \(t=u-v\), we obtain
\[
    \ind[u=v]
    =
    \frac{1}{q}\sum_{c\in\bbF_q}\omega_p^{\Tr(c(u-v))}.
\]
\end{proof}

\begin{definition}[Relative Hamming distance]
\label{def:relative-hamming-distance}
\lean{BlrPcp.distance}
\leanok
For $f,g:\bbF_q^n\to\bbF_q$, define
\[
    \delta(f,g)
    \coloneqq \Pr_{x\in\bbF_q^n}\left[f(x)\ne g(x)\right].
\]
\end{definition}

\section{Linear functions over finite fields}

\begin{definition}[Linear functions]
\label{def:linear-functions}
\lean{BlrPcp.linearFn, BlrPcp.IsLinear, BlrPcp.LinearSet}
\leanok
For $\alpha\in\bbF_q^n$, define the linear scalar function
$\ell_\alpha:\bbF_q^n\to\bbF_q$ by
\[
    \ell_\alpha(x) \coloneqq \inner{\alpha}{x}.
\]
A function $h:\bbF_q^n\to\bbF_q$ is linear if $h=\ell_\alpha$ for some
$\alpha\in\bbF_q^n$. The finite set of all linear scalar functions is
\[
    \lin
    \coloneqq
    \left\{\ell_\alpha \mid \alpha\in\bbF_q^n\right\}.
\]
\end{definition}

\begin{definition}[Distance to linearity]
\label{def:distance-to-linearity}
\lean{BlrPcp.distanceToLinear}
\uses{def:relative-hamming-distance,def:linear-functions}
\leanok
The distance from $f:\bbF_q^n\to\bbF_q$ to linearity is
\[
    \delta(f,\lin) \coloneqq \min_{g\in\lin}\delta(f,g).
\]
\end{definition}



\begin{lemma}[Well-separatedness of $\lin$]
\label{lem:linear-set-well-separated}
\lean{BlrPcp.linearSet_wellSeparated}
\uses{def:linear-functions,def:relative-hamming-distance}
\leanok
For any distinct $f,g \in \lin$, one has $\delta(f,g)= 1-1/|\bbF|$.
\end{lemma}
\begin{proof}\uses{def:linear-functions,def:relative-hamming-distance}\leanok
Let $a,b \in \bbF^n$ be the unique vectors such that $f(x) = \langle a, x\rangle$ and $g(x) = \langle b, x\rangle$ for each $x$.
Let $c = a-b$, which must be a non-zero vector. In particular, there exists $i \in [n]$ such that $c_i \neq 0$. Take any permutation $\sigma: [n] \rightarrow [n]$ such that $\sigma(i)=n$. Then
\begin{align*}
    \delta(f,g) & = 1- \Pr_{x \sim \bbF^n}[\langle c, x\rangle =0] \\
    & = 1 - \frac{1}{|\bbF|^{n-1}}\sum_{x_1,\dots,x_{n-1} \in \bbF}\frac{1}{|\bbF|}\sum_{x_n \in \bbF} \ind[c_i \cdot x_n = -\sum_{j \in [n]\setminus {i}} c_j \cdot x_{\sigma(j)}] \\
    & = 1 - \frac{1}{|\bbF|^{n-1}}\sum_{x_1,\dots,x_{n-1} \in \bbF}\frac{1}{|\bbF|}\sum_{x_n \in \bbF} \ind[x_n = -c_i^{-1} \cdot \sum_{j \in [n]\setminus {i}} c_j \cdot x_{\sigma(j)}] 
\end{align*}
where the last equality follows because $c_i \neq 0$. Now note that  $\ind[x_n = -c_i^{-1} \cdot \sum_{j \in [n]\setminus {i}} c_j \cdot x_{\sigma(j)}]$ evaluates to $1$ for exactly one value of $x_n$. Thus we continue from above and get
\begin{align*}
    \delta(f,g) = 1-\frac{1}{|\bbF|^{n-1}}\sum_{x_1,\dots,x_{n-1} \in \bbF}\frac{1}{|\bbF|} = 1-\frac{1}{|\bbF|} \, .
\end{align*}
\end{proof}


\begin{lemma}[Distance formula]
\label{lem:distance-formula}
\lean{BlrPcp.distance_formula_via_phase_fourier_coefficients}
\uses{def:relative-hamming-distance,def:fourier-expansion-field-valued,def:function-inner-product}
\leanok
Let $f,g:\bbF_q^n\to\bbF_q$.
Let $\{\fexp{\ff}{c}\}_{c\in\bbF_q^\times}$ and
$\{\fexp{\fg}{c}\}_{c\in\bbF_q^\times}$ be their Fourier expansions.
Then
\[
    \delta(f,g)
    = 1-\frac{1}{q}\left(1+
        \sum_{c\in\bbF_q^\times}
        \inner{\fexp{\ff}{c}}{\fexp{\fg}{c}}
    \right).
\]
\end{lemma}

\begin{proof}\uses{lem:character-sum-indicator,lem:fourier-plancherel}\leanok
By \cref{lem:character-sum-indicator},
\[
    \Pr_x[f(x)=g(x)]
    = \bbE_x\left[\frac{1}{q}\sum_{c\in\bbF_q}
        \omega_p^{\Tr(c(f(x)-g(x)))}\right].
\]
The term $c=0$ contributes $1/q$. For $c\ne 0$,
\[
    \omega_p^{\Tr(c(f(x)-g(x)))}
    = \fexp{\ff}{c}(x)\overline{\fexp{\fg}{c}(x)}.
\]
Therefore
\[
    \Pr_x[f(x)=g(x)]
    =\frac{1}{q}\left(1+\sum_{c\in\bbF_q^\times}
        \inner{\fexp{\ff}{c}}{\fexp{\fg}{c}}
    \right),
\]
which gives the first formula after using
$\delta(f,g)=1-\Pr_x[f(x)=g(x)]$.
\end{proof}

\begin{lemma}[Fourier expansion of a linear function]
\label{lem:linear-function-fourier-expansion}
\lean{BlrPcp.phaseLift_linearFn}
\uses{def:linear-functions,def:fourier-expansion-field-valued,def:additive-character}
\leanok
Fix $\alpha\in\bbF_q^n$ and let $\{\lambda_c\}_{c \in \bbF_q^\times}$ be the Fourier expansion of $\ell_\alpha$.
Then
\[
    \fexp{\lambda}{c}=\chi_{c\alpha}.
\]
\end{lemma}

\begin{proof}\leanok
For every $x\in\bbF_q^n$,
\[
    \fexp{\lambda}{c}(x)
    = \omega_p^{\Tr(c\inner{\alpha}{x})}
    = \omega_p^{\Tr(\inner{c\alpha}{x})}
    = \chi_{c\alpha}(x).
\]
\end{proof}

\begin{definition}[Linear Fourier score]
\label{def:linear-fourier-score}
\lean{BlrPcp.linearFourierScore, BlrPcp.linearFourierScoreC}
\uses{def:nonzero-field-elements,def:fourier-expansion-field-valued,def:fourier-coefficient}
\leanok
For $f:\bbF_q^n\to\bbF_q$ and $\alpha\in\bbF_q^n$, define the linear Fourier score
of $f$ at $\alpha$ by
\[
    S_\alpha(f) \coloneqq \sum_{c\in\bbF_q^\times}
        \widehat{\fexp{\ff}{c}}(c\alpha),
\]
where $\{\fexp{\ff}{c}\}_{c \in \bbF_q^\times}$ is the Fourier expansion of $f$.
\end{definition}

\begin{lemma}[Distance to linearity in Fourier form]
\label{lem:distance-to-linearity-fourier}
\lean{BlrPcp.distance_to_linearity_fourier}
\uses{def:distance-to-linearity,def:linear-fourier-score}
\leanok
Let $f:\bbF_q^n\to\bbF_q$ have Fourier expansion
$\{\fexp{\ff}{c}\}_{c\in\bbF_q^\times}$. Then
\[
    \delta(f,\lin)
    = 1-\frac{1}{q}\left(1+
        \max_{\alpha\in\bbF_q^n}
        S_\alpha(f)
    \right).
\]
Moreover, for every $\alpha\in\bbF_q^n$, the score $S_\alpha(f)$ is real and
satisfies $-1\le S_\alpha(f)\le q-1$.
\end{lemma}

\begin{proof}\uses{lem:distance-formula,lem:linear-function-fourier-expansion}\leanok
Apply \cref{lem:distance-formula} with $g=\ell_\alpha$ and use
\cref{lem:linear-function-fourier-expansion}. This gives
\[
    \delta(f,\ell_\alpha)
    =1-\frac{1}{q}\left(1+
        S_\alpha(f)
    \right).
\]
Taking the minimum over $\alpha$ gives the claimed formula for
$\delta(f,\lin)$.
For fixed $\alpha$, the same formula expresses $S_\alpha(f)$ as
$q(1-\delta(f,\ell_\alpha))-1$. Hence $S_\alpha(f)$ is real, and since
$0\le \delta(f,\ell_\alpha)\le 1$, we get $-1\le S_\alpha(f)\le q-1$.
\end{proof}

\section{BLR linearity test}

\begin{definition}[Finite-field BLR test]
\label{def:finite-field-blr-test}
\lean{BlrPcp.BLRAcceptsAt, BlrPcp.acceptanceProbabilityBLR}
\uses{def:nonzero-field-elements}
\leanok
Given oracle access to $f:\bbF_q^n\to\bbF_q$, the finite-field BLR verifier
$\BLRVerifier$ samples $a,b\in\bbF_q^\times$ and $x,y\in\bbF_q^n$ independently and uniformly at random, and accepts iff
\[
    a f(x)+b f(y)=f(ax+by).
\]
That is,
\[
    \BLRVerifier=1
    \quad\Longleftrightarrow\quad
    a f(x)+b f(y)=f(ax+by).
\]
\end{definition}

\begin{lemma}[Completeness]
\label{lem:blr-completeness}
\lean{BlrPcp.blr_completeness}
\uses{def:finite-field-blr-test,def:linear-functions}
\leanok
If $f\in\lin$, then
\[
    \Pr[\BLRVerifier=1]=1.
\]
\end{lemma}

\begin{proof}\leanok
If $f=\ell_\alpha$ for some $\alpha\in\bbF_q^n$, then for every
$a,b\in\bbF_q^\times$ and $x,y\in\bbF_q^n$,
\[
    f(ax+by)
    =\inner{\alpha}{ax+by}
    =a\inner{\alpha}{x}+b\inner{\alpha}{y}
    =a f(x)+b f(y).
\]
Thus the verifier accepts for every choice of randomness.
\end{proof}

\begin{lemma}[Exact BLR acceptance formula]
\label{lem:exact-blr-acceptance}
\lean{BlrPcp.exact_blr_acceptance_formula}
\uses{def:finite-field-blr-test,def:linear-fourier-score}
\leanok
Let $f:\bbF_q^n\to\bbF_q$.
Then
\[
    \Pr[\BLRVerifier=1]
    = \frac{1}{q}\left(1+\frac{1}{(q-1)^2}
        \sum_{\alpha\in\bbF_q^n} S_\alpha(f)^3
    \right).
\]
\end{lemma}

\begin{proof}\uses{lem:character-sum-indicator,lem:fourier-inversion,lem:character-orthonormal}\leanok
By \cref{lem:character-sum-indicator}, for fixed $a,b,x,y$,
\[
    \ind[a f(x)+b f(y)=f(ax+by)]
    =\frac{1}{q}\sum_{c\in\bbF_q}
        \omega_p^{\Tr(c(a f(x)+b f(y)-f(ax+by)))}.
\]
Taking expectation gives
\[
    \Pr[\BLRVerifier=1]
    =\frac{1}{q}+\frac{1}{q}\bbE_{a,b,x,y}
        \left[\sum_{c\in\bbF_q^\times}
        \fexp{\ff}{ca}(x)\fexp{\ff}{cb}(y)
        \fexp{\ff}{-c}(ax+by)
        \right].
\]
Expand the three complex-valued functions into their Fourier expansions:
\begin{align*}
&\bbE_{a,b,x,y}
        \left[\sum_{c\in\bbF_q^\times}
        \fexp{\ff}{ca}(x)\fexp{\ff}{cb}(y)
        \fexp{\ff}{-c}(ax+by)
        \right] \\
&=\bbE_{a,b}\left[\sum_{c\in\bbF_q^\times}
    \sum_{\alpha,\beta,\gamma\in\bbF_q^n}
    \widehat{\fexp{\ff}{ca}}(\alpha)
    \widehat{\fexp{\ff}{cb}}(\beta)
    \widehat{\fexp{\ff}{-c}}(\gamma)
    \bbE_{x,y}\left[
        \chi_\alpha(x)\chi_\beta(y)\chi_\gamma(ax+by)
    \right]
\right].
\end{align*}
Since
\[
    \chi_\gamma(ax+by)=\chi_{a\gamma}(x)\chi_{b\gamma}(y),
\]
\cref{lem:character-orthonormal} implies that the inner expectation is $1$ exactly when
\[
    \alpha+a\gamma=0
    \quad\text{and}\quad
    \beta+b\gamma=0,
\]
and is $0$ otherwise. Therefore the last display equals
\[
    \bbE_{a,b}\left[\sum_{c\in\bbF_q^\times}
    \sum_{\alpha\in\bbF_q^n}
    \widehat{\fexp{\ff}{ca}}(\alpha)
    \widehat{\fexp{\ff}{cb}}(a^{-1}b\alpha)
    \widehat{\fexp{\ff}{-c}}(-a^{-1}\alpha)
    \right].
\]
Substitute $\alpha=ca\eta$. Then this becomes
\[
    \bbE_{a,b}\left[\sum_{c\in\bbF_q^\times}
    \sum_{\eta\in\bbF_q^n}
    \widehat{\fexp{\ff}{ca}}(ca\eta)
    \widehat{\fexp{\ff}{cb}}(cb\eta)
    \widehat{\fexp{\ff}{-c}}(-c\eta)
    \right].
\]
As $a,b,c$ range over $\bbF_q^\times$, so do $ca$, $cb$, and $-c$. Hence this equals
\[
    \frac{1}{(q-1)^2}
    \sum_{\eta\in\bbF_q^n}
    \left(\sum_{d\in\bbF_q^\times}
        \widehat{\fexp{\ff}{d}}(d\eta)
    \right)^3
    =\frac{1}{(q-1)^2}\sum_{\eta\in\bbF_q^n}S_\eta(f)^3.
\]
Substituting into the expression for $\Pr[\BLRVerifier=1]$ proves the claim.
\end{proof}

\begin{lemma}[Second moment bound]
\label{lem:second-moment-bound}
\lean{BlrPcp.sum_linearFourierScore_sq_upperbound}
\uses{def:linear-fourier-score}
\leanok
For every $f:\bbF_q^n\to\bbF_q$, one has
\[
    \sum_{\alpha\in\bbF_q^n} S_\alpha(f)^2 \le (q-1)^2.
\]
\end{lemma}

\begin{proof}\uses{lem:parseval-identity}\leanok
Expanding the square gives
\[
    \sum_{\alpha\in\bbF_q^n} S_\alpha(f)^2
    =\sum_{a,b\in\bbF_q^\times}
        \sum_{\alpha\in\bbF_q^n}
        \widehat{\fexp{\ff}{a}}(a\alpha)
        \widehat{\fexp{\ff}{b}}(b\alpha).
\]
For fixed $a,b\in\bbF_q^\times$, Cauchy-Schwarz and
\cref{lem:parseval-identity} give
\begin{align*}
\abs{\sum_{\alpha\in\bbF_q^n}
        \widehat{\fexp{\ff}{a}}(a\alpha)
        \widehat{\fexp{\ff}{b}}(b\alpha)}
&\le
\left(\sum_{\alpha\in\bbF_q^n}\abs{\widehat{\fexp{\ff}{a}}(a\alpha)}^2\right)^{1/2}
\left(\sum_{\alpha\in\bbF_q^n}\abs{\widehat{\fexp{\ff}{b}}(b\alpha)}^2\right)^{1/2} \\
&=1.
\end{align*}
The equality uses that multiplication by $a$ and by $b$ permutes $\bbF_q^n$, and that
Parseval plus $\abs{\fexp{\ff}{a}(x)}=\abs{\fexp{\ff}{b}(x)}=1$ for every $x$
make both squared $L_2$ norms equal to \(1\).
Summing over the $(q-1)^2$ choices of $(a,b)$ gives the result.
\end{proof}

\begin{theorem}[Soundness of the finite-field BLR test]
\label{thm:finite-field-blr-soundness}
\lean{BlrPcp.blr_soundness}
\uses{def:finite-field-blr-test,def:distance-to-linearity}
\leanok
For every function $f:\bbF_q^n\to\bbF_q$, one has
\[
    \Pr[\BLRVerifier=1] \le 1-\delta(f,\lin).
\]
\end{theorem}

\begin{proof}\uses{lem:distance-to-linearity-fourier,lem:exact-blr-acceptance,lem:second-moment-bound}\leanok
Let $S_\alpha(f)$ be as in \cref{def:linear-fourier-score}, and let
\[
    M \coloneqq \max_{\alpha\in\bbF_q^n}S_\alpha(f).
\]
By \cref{lem:distance-to-linearity-fourier},
\[
    1-\delta(f,\lin)=\frac{1}{q}(1+M).
\]
By \cref{lem:distance-to-linearity-fourier}, each $S_\alpha(f)$ is real and
$S_\alpha(f)\le M$.
Therefore $S_\alpha(f)^3\le M S_\alpha(f)^2$ for every $\alpha$, because
$S_\alpha(f)^2\ge 0$.
Using \cref{lem:exact-blr-acceptance} and \cref{lem:second-moment-bound},
\begin{align*}
    \Pr[\BLRVerifier=1]
    &=\frac{1}{q}\left(1+\frac{1}{(q-1)^2}
        \sum_{\alpha\in\bbF_q^n}S_\alpha(f)^3\right) \\
    &\le \frac{1}{q}\left(1+\frac{M}{(q-1)^2}
        \sum_{\alpha\in\bbF_q^n}S_\alpha(f)^2\right) \\
    &\le \frac{1}{q}(1+M) \\
    &=1-\delta(f,\lin).
\end{align*}
\end{proof}

% Below is the GH-route to the same soundness results.
\begin{proof}[Alternative proof via Gowers–Hatami]\uses{lem:rep-iff-blr, lem:gh_linearity_formalized}
First denote
\[
X := \F_q^n, \qquad A := \F_q^\times, \qquad \omega := e^{2\pi i/p}.
\]
Define a group $G$ on the set $X \times A$ with the binary operation
\[
(x, a) \star (y, b) := \bigl(b^{-1}x + a^{-1}y,\; ab\bigr).
\]
Let $F_f$ be the unitary matrix valued lift $F_f : G \to U(\cc^{q-1})$ given by the diagonal matrix
\[
F_f(x, a) = D(a f(x)) := \sum_{c \in A} \omega^{\Tr\, c a f(x)}\ket{c}\bra{c},
\]
where $\{\ket{c}\}$ denotes the orthonormal basis vectors of a Hilbert space $\mathcal{H} \cong \cc^{q-1}$, the
inner product between operators is $\inner{A}{B}_\sigma=\Tr(A^* B \sigma)$, and $\sigma = \I/(q-1)$. By
\lemref{lem:rep-iff-blr}, since $f$ is rejected with probability $\varepsilon$, the lift $F_f$ is a
$(\frac{q}{q-1}\varepsilon,\frac{\I}{q-1})$ approximate representation of $(G, \star)$.

Then by \lemref{lem:gh_linearity_formalized}, there exists a linear function $\ell: \FF_q^n \to \FF_q$ whose
relative Hamming distance to $f$ is bounded by
\[
    \delta(f, \ell) \le \frac{q-1}{q}\cdot\frac{q}{q-1}\varepsilon = \varepsilon.
\]
This proves the claim
\[
\delta(f, \lin) := \min_{\alpha} \delta(f, \ell_\alpha) \leq \delta(f, \ell) \leq \varepsilon.
\]
Miraculously, the loose factor $\frac{q}{q-1}$ in the approximate-representation parameter cancels the tighter
factor $\frac{q-1}{q}$ of \lemref{lem:gh_linearity_formalized}.

In Lean this route is formalized as \texttt{gh\_blr\_soundness\_acceptance}, which proves
\texttt{acceptanceProbabilityBLR f $\le$ 1 - distanceToLinear f} (the statement of \texttt{blr\_soundness}),
valid for $\frac{q}{q-1}\varepsilon\le1$; the unconditional statement is the direct proof above, since for
$\frac{q}{q-1}\varepsilon>1$ the third-moment extraction's $\max\ge0$ step fails.
\end{proof}
